%% Sample Exam 7 -- Medical Image Analysis (ETH Zurich, 2026)
%%
%% Written in the style of Example_Exam_with_Answers.pdf.
%%
%% Compile the version WITH solutions:      pdflatex "\def\withanswers{}%% Sample Exam 7 -- Medical Image Analysis (ETH Zurich, 2026)
%%
%% Written in the style of Example_Exam_with_Answers.pdf.
%%
%% Compile the version WITH solutions:      pdflatex "\def\withanswers{}%% Sample Exam 7 -- Medical Image Analysis (ETH Zurich, 2026)
%%
%% Written in the style of Example_Exam_with_Answers.pdf.
%%
%% Compile the version WITH solutions:      pdflatex "\def\withanswers{}%% Sample Exam 7 -- Medical Image Analysis (ETH Zurich, 2026)
%%
%% Written in the style of Example_Exam_with_Answers.pdf.
%%
%% Compile the version WITH solutions:      pdflatex "\def\withanswers{}\input{Sample_Exam_7.tex}"
%% Compile the blank version (no answers):  pdflatex Sample_Exam_7.tex
%%
\ifx\withanswers\undefined
  \documentclass[11pt]{exam}
\else
  \documentclass[11pt,answers]{exam}
\fi

\usepackage[margin=2.4cm]{geometry}
\usepackage{amsmath,amssymb}
\usepackage{array}
\usepackage[T1]{fontenc}

% ---------------------------------------------------------------- page style
\pagestyle{foot}
\firstpageheader{}{}{}
\firstpagefooter{}{}{}
\runningheader{}{}{}
\runningfooter{}{}{Page \thepage}

% -------------------------------------------------------------- multiplechoice
% Reproduce the "(A) (B) (C) ..." labelling of the original exam.
\renewcommand{\choicelabel}{$\bigcirc$~(\Alph{choice})}

% ---------------------------------------------------------- true/false macros
% \TFT -> the statement is TRUE, \TFF -> the statement is FALSE.
% In the answer version the correct word is set in capitals and bold.
\newcommand{\TFT}{\ifprintanswers\textbf{TRUE}\hspace{1.1em}False\else True\hspace{1.1em}False\fi}
\newcommand{\TFF}{True\hspace{1.1em}\ifprintanswers\textbf{FALSE}\else False\fi}

% ------------------------------------------------------------- writing space
% Reserves room to write in the blank version; collapses in the answer version.
\newcommand{\answerspace}[1]{\ifprintanswers\else\vspace{#1}\fi}

\begin{document}

% ------------------------------------------------------------------ title page
\begin{center}
  {\LARGE\bfseries Sample Exam}\\[2mm]
  {\Large Medical Image Analysis}\\[6mm]
\end{center}

\vspace{4mm}
\noindent Student name:\hrulefill

\vspace{6mm}
\noindent Immatriculation number:\hrulefill

\vspace{10mm}

\begin{center}
  {\large\bfseries Exam Questions}
\end{center}

\noindent\textbf{Note:} All the questions have equal weight.

\vspace{4mm}

\begin{questions}

% =============================================================== Question 1 ===
\question
In 2D acquisitions, slices are acquired in one pre-defined plane and stacked along the remaining
direction: axial slices stack head-to-toe, sagittal slices left-to-right, coronal slices
anterior-to-posterior. Voxel sizes are given as in-plane (two values) and through-plane (the slice
thickness), and the field-of-view is image size $\times$ voxel size.

\vspace{1mm}
\noindent You are handed a volume with the following properties, but the acquisition plane is not
recorded: field-of-view $240$~mm head-to-toe, $240$~mm left-to-right, $180$~mm
anterior-to-posterior; voxel size in-plane $0.75 \times 0.75$~mm$^2$, through-plane $4.5$~mm.

\begin{parts}

  \part In which plane was this volume acquired? Justify your answer from the numbers.

  \begin{solution}
    Answer: a \textbf{coronal} acquisition.

    The through-plane direction is the one whose extent is sampled at $4.5$~mm, and only
    anterior-to-posterior works: $180 / 4.5 = 40$ slices, an integer. If either of the other two
    directions were through-plane, it would need $240 / 4.5 = 53.3$ slices --- not an integer, so
    the numbers are inconsistent with any other assignment. Slices stacked along
    anterior-to-posterior lie in the coronal plane.
  \end{solution}
\answerspace{3.0cm}

  \part How many voxels does the volume have in the head-to-toe, left-to-right and
  anterior-to-posterior directions, respectively?

  \begin{solution}
    Answer: $320 \times 320 \times 40$~voxels.

    Head-to-toe and left-to-right are in-plane: $240 / 0.75 = 320$ each. Anterior-to-posterior is
    through-plane: $180 / 4.5 = 40$.
  \end{solution}
\answerspace{2.2cm}

  \part The volume is resampled to isotropic $1.5 \times 1.5 \times 1.5$~mm$^3$ voxels. Give the
  resize factors and the resulting image size, and verify that the field-of-view is unchanged.

  \begin{solution}
    Answer: resize factors $0.5 \times 0.5 \times 3.0$, resulting image size
    $160 \times 160 \times 120$~voxels.

    The factor is the ratio of old to new voxel size: $0.75/1.5 = 0.5$ in-plane,
    $4.5/1.5 = 3.0$ through-plane. Sizes: $320 \cdot 0.5 = 160$, $320 \cdot 0.5 = 160$,
    $40 \cdot 3.0 = 120$. Field-of-view check: $160 \cdot 1.5 = 240$, $160 \cdot 1.5 = 240$,
    $120 \cdot 1.5 = 180$~mm --- unchanged, as it must be. (The anterior-to-posterior direction is
    upsampled by 3: two of every three samples there are interpolated, not measured.)
  \end{solution}
\answerspace{3.6cm}

\end{parts}

% =============================================================== Question 2 ===
\question
In bilinear interpolation the intensity at a query point $x_1$ (with $x_1^0 \le x_1 \le x_1^1$) is
first interpolated along one dimension from the two neighbouring pixels. Which of the two
expressions below is the correct linear interpolation? Please circle the correct choice.

\vspace{2mm}
\begin{center}
  $\displaystyle I(x_1) = \frac{x_1^1 - x_1}{x_1^1 - x_1^0}\, I(x_1^0)
  + \frac{x_1 - x_1^0}{x_1^1 - x_1^0}\, I(x_1^1)$
  \hspace{1.1cm}
  $\displaystyle I(x_1) = \frac{x_1 - x_1^0}{x_1^1 - x_1^0}\, I(x_1^0)
  + \frac{x_1^1 - x_1}{x_1^1 - x_1^0}\, I(x_1^1)$
\end{center}
\vspace{2mm}

\begin{solution}
  Answer: Left one.

  Each pixel's weight must be proportional to the query point's distance from the \emph{other}
  pixel --- nearness earns weight. The sanity check takes two seconds: set $x_1 = x_1^0$. The left
  expression gives $1 \cdot I(x_1^0) + 0 \cdot I(x_1^1) = I(x_1^0)$, as an interpolant must. The
  right expression gives $I(x_1^1)$ --- the wrong pixel entirely; its weights are inverted.
  Bilinear interpolation applies this step along $x_1$ and then along $x_2$ (trilinear adds the
  third dimension, using the 8 cube corners).
\end{solution}
\answerspace{1.5cm}

% =============================================================== Question 3 ===
\question
Probabilistic reasoning underlies most of this course.

\begin{parts}
  \part If each pixel intensity is treated as a realization of a random variable, the (normalized)
  histogram and the cumulative histogram of an image approximate two standard probabilistic
  objects. Which two --- and what assumption about the pixels makes this approximation valid?

  \begin{solution}
    Answer: the normalized histogram approximates the \textbf{probability density function} of the
    intensity, and the cumulative histogram approximates the \textbf{cumulative distribution
    function}.

    The assumption is that each pixel intensity is an \textbf{independent realization} of the same
    random variable $I$ --- only then is counting pixels per intensity bin an estimate of the
    underlying distribution.
  \end{solution}
\answerspace{2.8cm}

  \part The Maximum-A-Posteriori estimate is $\arg\max_l\, p(i \mid l)\, p(l)$ and the Maximum
  Likelihood estimate is $\arg\max_l\, p(i \mid l)$. Under what condition on the prior are the two
  identical?

  \begin{solution}
    Answer: when the prior is \textbf{uniform}, $p(l) = c$ for all $l$.

    A constant factor does not change the location of a maximum, so the prior drops out of the
    $\arg\max$ and MAP reduces to MLE. Conversely, whenever anyone uses MLE they are implicitly
    assuming that no value of $l$ is more plausible than any other before seeing the data.
  \end{solution}
\answerspace{2.2cm}
\end{parts}

% =============================================================== Question 4 ===
\question
Non-local means computes $\hat I(x) = \sum_{y} w(x,y) J(y)$ with weights from the similarity of
the patches $N(x)$ and $N(y)$, a degree-of-filtering parameter $h$, and averaging restricted to a
search window $W(x)$. Please indicate whether the statement is true or false by circling the
correct choice in each of the following.

\begin{parts}
  \part \TFF\ \ Choosing the patch to be the entire image, $N(x) = \Omega$, reduces non-local means
  to Gaussian filtering.

  \part \TFT\ \ Choosing the patch to be the single pixel, $N(x) = \{x\}$, reduces non-local means
  to Yaroslavsky's filter, a special case of bilateral filtering.

  \part \TFF\ \ Increasing $h$ decreases the amount of smoothing.

  \part \TFT\ \ The averaging is restricted to a search window $W(x)$ rather than the whole image
  mainly to reduce the computational cost.
\end{parts}

\begin{solution}
  (a) FALSE --- it reduces to \emph{mean} filtering, not Gaussian filtering. If every pixel's patch
  is the whole image, all patches are identical, every weight is equal, and the weighted average
  degenerates to the plain average.

  (b) TRUE --- with a single-pixel ``patch'' the weight compares raw grey values, which is exactly
  the neighbourhood-filter family (Yaroslavsky, bilateral). This is also why NLM is more robust:
  a single noisy value is a fragile basis for similarity, a whole patch is not.

  (c) FALSE --- $h$ is the degree of filtering: larger $h$ makes the exponential weights decay more
  slowly, so more (less similar) patches contribute and the result is \emph{smoother}. $h$ should
  be tied to the noise level, which is worth estimating first.

  (d) TRUE --- and it also encodes a mild prior that useful matches are nearby. The three
  parameters of the method are exactly $h$, the patch size, and the search-window size.
\end{solution}

% =============================================================== Question 5 ===
\question
A 2D landmark-based registration uses thin-plate splines with the $N = 50$ landmarks as control
points, combined with a global affine part. How many parameters does the transformation have? How
does this compare with a naive deformable parameterization of a $256 \times 256$ image that stores
one independent displacement vector per pixel?

\begin{solution}
  Answer: the TPS model has $50 \times 2 + 6 = \mathbf{106}$ parameters; the naive field has
  $256 \times 256 \times 2 = \mathbf{131\,072}$ --- about $1\,240$ times more.

  Each control point carries a coefficient vector $d_n \in \mathbb{R}^2$ (100 parameters), and the
  2D affine part adds $4 + 2 = 6$. The comparison is the point of the question: both are
  ``deformable'' models, but the TPS parameterization builds smoothness in and needs three orders
  of magnitude fewer parameters, while the naive field is not only huge but can produce noisy,
  topology-breaking transformations for generic parameter values.
\end{solution}
\answerspace{3.6cm}

% =============================================================== Question 6 ===
\question
An Active Appearance Model extends a statistical shape model with a model of texture. Please
describe briefly (i) how the texture model is built from the training images, and (ii) how shape
and texture are combined into a single set of appearance parameters. You do not need to write down
the equations.

\begin{solution}
  Answer: \emph{(i) Building the texture model.} Every training image is \textbf{warped to the mean
  shape}, which removes the shape variation and yields a shape-free patch. The intensities are
  sampled from this shape-normalized image into a texture vector and \textbf{normalized in
  intensity} to eliminate lighting variation. PCA is then applied to the normalized texture
  vectors, exactly as it was applied to the shape coordinates: $g \approx \bar g + P_g b_g$.

  \emph{(ii) Combining shape and texture.} Each training sample is now summarized by two parameter
  vectors, $b_s$ for shape and $b_g$ for texture --- and the two are correlated (a shape change
  usually comes with a texture change). They are stacked into a single vector, with a weighting
  $W_s$ on the shape part to account for the different units of shape and texture, and a
  \textbf{second PCA} is run on the stacked vectors. The resulting parameters $c$ control shape and
  texture \emph{jointly}: a new image is generated by producing the shape-free texture from $c$ and
  warping it with the shape that the same $c$ describes.

  (Fitting then works by synthesising the model image and minimising its difference to the target,
  with the parameter-update matrix precomputed from training-set perturbations.)
\end{solution}
\answerspace{7.5cm}

% =============================================================== Question 7 ===
\question
Which of the following statements about atlas-based segmentation practice is \textbf{inaccurate}?

\begin{choices}
  \choice An atlas consists of an image together with its manual segmentation; it is registered to
  the unseen patient image and the labels are propagated along the transformation.
  \choice When an average (mean) atlas is built from a population, the choice of the reference
  subject does not influence the result --- averages built from different references are the same.
  \choice The quality of an atlas-based segmentation depends on the image similarity between atlas
  and patient, the presence of pathology, the quality of the registration, and how typical the
  patient is of the atlas population.
  \choice Shape-based averaging fuses several segmentations by averaging their distance transforms
  rather than voting on their labels.
\end{choices}

\begin{solution}
  Answer: B.

  This is the \textbf{reference-selection bias}: average models built with different reference
  subjects are demonstrably \emph{not} the same --- the lecture shows two different average images
  from two different references. The bias is precisely what motivates unbiased, group-wise
  template construction in a natural coordinate frame, rather than anchoring the average to any
  single individual.
\end{solution}

% =============================================================== Question 8 ===
\question
Which of the following statements on uncertainty estimation is \textbf{incorrect}?

\begin{choices}
  \choice One source of uncertainty is that the input $x$ may not uniquely identify the label $y$.
  \choice Another source is that the training set $D$ is itself a random sample from the
  population, and a different draw would have produced different parameters.
  \choice The Generalized Energy Distance can be evaluated on an ordinary test set with a single
  annotation per image, which is why it is the most general evaluation of uncertainty.
  \choice The Expected Calibration Error measures the average gap between the confidence a model
  assigns and the accuracy it actually achieves at that confidence.
\end{choices}

\begin{solution}
  Answer: C.

  The two evaluation routes are exactly reversed. GED (like MMD) is a \emph{distributional}
  distance between the set of model samples and the set of \textbf{observed annotations} --- it
  needs data sets with \emph{multiple labels per image}, such as multi-rater segmentations. When
  only one annotation per image is available --- the most general case --- the tool is
  \textbf{calibration}, i.e.\ ECE: whether the assigned confidence tracks the achieved accuracy.
\end{solution}

% =============================================================== Question 9 ===
\question
Which of the following statements from the machine-learning recap is \textbf{incorrect}?

\begin{choices}
  \choice Hyper-parameters --- the architecture, the stopping point of the optimization, the
  learning rate --- are chosen using a validation set, ideally as a nested optimization inside the
  training of the parameters.
  \choice The training loss is a sum over samples, so the optimal parameters do not have to satisfy
  any single sample perfectly --- they satisfy the samples on average.
  \choice The prediction error on the test set must be computed with the same differentiable loss
  that was used for training, since otherwise the numbers are not comparable.
  \choice For a new sample the prediction is $y \approx f(x; \theta^*)$, and in general
  $y \neq f(x; \theta^*)$.
\end{choices}

\begin{solution}
  Answer: C.

  Backwards on both counts. Differentiability is a requirement of \emph{training} --- the loss must
  be optimized --- not of \emph{evaluation}. On the test set one may, and often should, use a
  different, more application-appropriate function $D$, precisely in the cases where $D$ is not
  differentiable (a ranking-based measure is the lecture's example). The same logic appeared with
  segmentation: the set-based Dice score is the standard \emph{evaluation} measure even when the
  network was trained with cross-entropy or soft Dice.
\end{solution}

% ============================================================== Question 10 ===
\question
``Meaningful perturbation'' explains a classifier's decision by perturbing the input. Which of the
following statements describes the method correctly?

\begin{choices}
  \choice It searches the training set for the image that most strongly activates the class and
  presents it as the explanation.
  \choice It optimizes an information-deletion mask --- applying blur, noise or masking --- that
  covers as little of the image as possible while dropping the class score as much as possible.
  \choice It requires no evaluations of the model, which makes it the cheapest explanation method.
  \choice It computes the gradient of the class score with respect to the convolutional feature
  maps and pools it into a heatmap.
  \choice Because the mask is optimized, the resulting perturbations are guaranteed to remain
  natural-looking.
\end{choices}

\begin{solution}
  Answer: B.

  The key idea is to find the \emph{right} perturbation --- a form of information deletion --- and
  study its effect on the prediction: a mask objective balances minimal coverage against maximal
  drop in the class score. (D) describes Grad-CAM, a different method. (E) is the method's known
  caveat, inverted: unconstrained optimization can produce \emph{unnatural} perturbations ---
  adversarial-like artifacts against the explanation --- which is why the meaning of the
  explanation depends on the meaning of the changes applied to the input.
\end{solution}

% ============================================================== Question 11 ===
\question
A classification network takes grey-scale input images of size $68 \times 68$ with one input
channel. All convolutions are \emph{valid} (no padding) with stride $S = 1$:

\begin{itemize}
  \item \textbf{Conv1}: 6 kernels with $K = 5 \times 5$ \quad\textbf{Pool1}: max pooling
    $2 \times 2$, $S = 2$
  \item \textbf{Conv2}: 12 kernels with $K = 3 \times 3$ \quad\textbf{Pool2}: max pooling
    $2 \times 2$, $S = 2$
  \item \textbf{Conv3}: 24 kernels with $K = 3 \times 3$
\end{itemize}

\noindent Circle True or False:

\begin{parts}
  \part \TFF\ \ The output of Conv3 has a spatial size of $15 \times 15$.

  \part \TFT\ \ The global receptive field of a neuron in Conv3 is $20 \times 20$.

  \part \TFF\ \ The global stride of a neuron in Conv3 is $8 \times 8$.

  \part \TFT\ \ The network's convolutional layers have
  $5 \cdot 5 \cdot 1 \cdot 6 + 3 \cdot 3 \cdot 6 \cdot 12 + 3 \cdot 3 \cdot 12 \cdot 24 = 3\,390$
  kernel weights, and including one bias per kernel the total number of parameters is $3\,432$.
\end{parts}

\begin{solution}
  (a) FALSE --- $15 \times 15$ is the size \emph{entering} Conv3, not leaving it. Tracking:
  $68 \xrightarrow{\text{Conv1}} 64 \xrightarrow{\text{Pool1}} 32 \xrightarrow{\text{Conv2}} 30
  \xrightarrow{\text{Pool2}} 15 \xrightarrow{\text{Conv3}} 13$. The output is $13 \times 13$.

  (b) TRUE --- track the receptive field $R$ and jump $j$ from $R = 1$, $j = 1$ with
  $R \leftarrow R + (K-1)\cdot j$ and $j \leftarrow j \cdot S$:
  Conv1 $\rightarrow R = 5,\ j = 1$; Pool1 $\rightarrow R = 6,\ j = 2$;
  Conv2 $\rightarrow R = 10,\ j = 2$; Pool2 $\rightarrow R = 12,\ j = 4$;
  Conv3 $\rightarrow R = 20,\ j = 4$.

  (c) FALSE --- there are only \emph{two} stride-2 poolings, so the global stride is
  $2 \cdot 2 = 4$, i.e.\ $4 \times 4$. Strides compound multiplicatively over the layers that
  actually stride; Conv3 itself has stride 1 and changes nothing.

  (d) TRUE --- weights: $150 + 648 + 2\,592 = 3\,390$. Biases: one per kernel, i.e.\ per output
  channel: $6 + 12 + 24 = 42$, giving $3\,432$ in total. Pooling layers contribute neither weights
  nor biases.
\end{solution}

% ============================================================== Question 12 ===
\question
Which of the following statements are correct (multiple answers are possible)?

\begin{choices}
  \choice In FCN-style segmentation, outputs taken at several depths of the network are upsampled
  by different factors and combined, so that fine detail from shallow layers complements the
  context of deep layers.
  \choice A label map may be transformed with bilinear interpolation provided the interpolated
  values are rounded to the nearest integer label afterwards.
  \choice Matching the mean and standard deviation of two images,
  $f(j) = (j - \mu_j)\,\sigma_i/\sigma_j + \mu_i$, matches second-order statistics and is less
  sensitive to outliers than min--max normalization.
  \choice A similarity transformation has 4 degrees of freedom in 2D and 7 in 3D.
  \choice The curl of a transformation field, $\nabla \times T$, quantifies the local volumetric
  change induced by the transformation.
\end{choices}

\begin{solution}
  Answer: (A), (C) and (D) are correct.

  (B) is wrong, and the proposed repair makes it worse rather than better. Bilinear interpolation
  between labels produces values that are not labels, and rounding them does not restore meaning:
  between a voxel of label 1 and a voxel of label 3, the interpolated value $2$ rounds to label 2
  --- possibly a completely different structure that is not even adjacent to either. Label maps are
  transformed with nearest-neighbour interpolation, full stop.

  (E) attributes the wrong quantity: the \textbf{curl} carries information about the amount of
  \emph{infinitesimal rotation}. Local volumetric change is quantified by the \textbf{Jacobian
  determinant} ($\det J > 1$ expansion, $< 1$ compression, $= 1$ none), with the divergence ---
  the trace of the Jacobian --- carrying related compression/expansion information.
\end{solution}

\end{questions}

\end{document}
"
%% Compile the blank version (no answers):  pdflatex Sample_Exam_7.tex
%%
\ifx\withanswers\undefined
  \documentclass[11pt]{exam}
\else
  \documentclass[11pt,answers]{exam}
\fi

\usepackage[margin=2.4cm]{geometry}
\usepackage{amsmath,amssymb}
\usepackage{array}
\usepackage[T1]{fontenc}

% ---------------------------------------------------------------- page style
\pagestyle{foot}
\firstpageheader{}{}{}
\firstpagefooter{}{}{}
\runningheader{}{}{}
\runningfooter{}{}{Page \thepage}

% -------------------------------------------------------------- multiplechoice
% Reproduce the "(A) (B) (C) ..." labelling of the original exam.
\renewcommand{\choicelabel}{$\bigcirc$~(\Alph{choice})}

% ---------------------------------------------------------- true/false macros
% \TFT -> the statement is TRUE, \TFF -> the statement is FALSE.
% In the answer version the correct word is set in capitals and bold.
\newcommand{\TFT}{\ifprintanswers\textbf{TRUE}\hspace{1.1em}False\else True\hspace{1.1em}False\fi}
\newcommand{\TFF}{True\hspace{1.1em}\ifprintanswers\textbf{FALSE}\else False\fi}

% ------------------------------------------------------------- writing space
% Reserves room to write in the blank version; collapses in the answer version.
\newcommand{\answerspace}[1]{\ifprintanswers\else\vspace{#1}\fi}

\begin{document}

% ------------------------------------------------------------------ title page
\begin{center}
  {\LARGE\bfseries Sample Exam}\\[2mm]
  {\Large Medical Image Analysis}\\[6mm]
\end{center}

\vspace{4mm}
\noindent Student name:\hrulefill

\vspace{6mm}
\noindent Immatriculation number:\hrulefill

\vspace{10mm}

\begin{center}
  {\large\bfseries Exam Questions}
\end{center}

\noindent\textbf{Note:} All the questions have equal weight.

\vspace{4mm}

\begin{questions}

% =============================================================== Question 1 ===
\question
In 2D acquisitions, slices are acquired in one pre-defined plane and stacked along the remaining
direction: axial slices stack head-to-toe, sagittal slices left-to-right, coronal slices
anterior-to-posterior. Voxel sizes are given as in-plane (two values) and through-plane (the slice
thickness), and the field-of-view is image size $\times$ voxel size.

\vspace{1mm}
\noindent You are handed a volume with the following properties, but the acquisition plane is not
recorded: field-of-view $240$~mm head-to-toe, $240$~mm left-to-right, $180$~mm
anterior-to-posterior; voxel size in-plane $0.75 \times 0.75$~mm$^2$, through-plane $4.5$~mm.

\begin{parts}

  \part In which plane was this volume acquired? Justify your answer from the numbers.

  \begin{solution}
    Answer: a \textbf{coronal} acquisition.

    The through-plane direction is the one whose extent is sampled at $4.5$~mm, and only
    anterior-to-posterior works: $180 / 4.5 = 40$ slices, an integer. If either of the other two
    directions were through-plane, it would need $240 / 4.5 = 53.3$ slices --- not an integer, so
    the numbers are inconsistent with any other assignment. Slices stacked along
    anterior-to-posterior lie in the coronal plane.
  \end{solution}
\answerspace{3.0cm}

  \part How many voxels does the volume have in the head-to-toe, left-to-right and
  anterior-to-posterior directions, respectively?

  \begin{solution}
    Answer: $320 \times 320 \times 40$~voxels.

    Head-to-toe and left-to-right are in-plane: $240 / 0.75 = 320$ each. Anterior-to-posterior is
    through-plane: $180 / 4.5 = 40$.
  \end{solution}
\answerspace{2.2cm}

  \part The volume is resampled to isotropic $1.5 \times 1.5 \times 1.5$~mm$^3$ voxels. Give the
  resize factors and the resulting image size, and verify that the field-of-view is unchanged.

  \begin{solution}
    Answer: resize factors $0.5 \times 0.5 \times 3.0$, resulting image size
    $160 \times 160 \times 120$~voxels.

    The factor is the ratio of old to new voxel size: $0.75/1.5 = 0.5$ in-plane,
    $4.5/1.5 = 3.0$ through-plane. Sizes: $320 \cdot 0.5 = 160$, $320 \cdot 0.5 = 160$,
    $40 \cdot 3.0 = 120$. Field-of-view check: $160 \cdot 1.5 = 240$, $160 \cdot 1.5 = 240$,
    $120 \cdot 1.5 = 180$~mm --- unchanged, as it must be. (The anterior-to-posterior direction is
    upsampled by 3: two of every three samples there are interpolated, not measured.)
  \end{solution}
\answerspace{3.6cm}

\end{parts}

% =============================================================== Question 2 ===
\question
In bilinear interpolation the intensity at a query point $x_1$ (with $x_1^0 \le x_1 \le x_1^1$) is
first interpolated along one dimension from the two neighbouring pixels. Which of the two
expressions below is the correct linear interpolation? Please circle the correct choice.

\vspace{2mm}
\begin{center}
  $\displaystyle I(x_1) = \frac{x_1^1 - x_1}{x_1^1 - x_1^0}\, I(x_1^0)
  + \frac{x_1 - x_1^0}{x_1^1 - x_1^0}\, I(x_1^1)$
  \hspace{1.1cm}
  $\displaystyle I(x_1) = \frac{x_1 - x_1^0}{x_1^1 - x_1^0}\, I(x_1^0)
  + \frac{x_1^1 - x_1}{x_1^1 - x_1^0}\, I(x_1^1)$
\end{center}
\vspace{2mm}

\begin{solution}
  Answer: Left one.

  Each pixel's weight must be proportional to the query point's distance from the \emph{other}
  pixel --- nearness earns weight. The sanity check takes two seconds: set $x_1 = x_1^0$. The left
  expression gives $1 \cdot I(x_1^0) + 0 \cdot I(x_1^1) = I(x_1^0)$, as an interpolant must. The
  right expression gives $I(x_1^1)$ --- the wrong pixel entirely; its weights are inverted.
  Bilinear interpolation applies this step along $x_1$ and then along $x_2$ (trilinear adds the
  third dimension, using the 8 cube corners).
\end{solution}
\answerspace{1.5cm}

% =============================================================== Question 3 ===
\question
Probabilistic reasoning underlies most of this course.

\begin{parts}
  \part If each pixel intensity is treated as a realization of a random variable, the (normalized)
  histogram and the cumulative histogram of an image approximate two standard probabilistic
  objects. Which two --- and what assumption about the pixels makes this approximation valid?

  \begin{solution}
    Answer: the normalized histogram approximates the \textbf{probability density function} of the
    intensity, and the cumulative histogram approximates the \textbf{cumulative distribution
    function}.

    The assumption is that each pixel intensity is an \textbf{independent realization} of the same
    random variable $I$ --- only then is counting pixels per intensity bin an estimate of the
    underlying distribution.
  \end{solution}
\answerspace{2.8cm}

  \part The Maximum-A-Posteriori estimate is $\arg\max_l\, p(i \mid l)\, p(l)$ and the Maximum
  Likelihood estimate is $\arg\max_l\, p(i \mid l)$. Under what condition on the prior are the two
  identical?

  \begin{solution}
    Answer: when the prior is \textbf{uniform}, $p(l) = c$ for all $l$.

    A constant factor does not change the location of a maximum, so the prior drops out of the
    $\arg\max$ and MAP reduces to MLE. Conversely, whenever anyone uses MLE they are implicitly
    assuming that no value of $l$ is more plausible than any other before seeing the data.
  \end{solution}
\answerspace{2.2cm}
\end{parts}

% =============================================================== Question 4 ===
\question
Non-local means computes $\hat I(x) = \sum_{y} w(x,y) J(y)$ with weights from the similarity of
the patches $N(x)$ and $N(y)$, a degree-of-filtering parameter $h$, and averaging restricted to a
search window $W(x)$. Please indicate whether the statement is true or false by circling the
correct choice in each of the following.

\begin{parts}
  \part \TFF\ \ Choosing the patch to be the entire image, $N(x) = \Omega$, reduces non-local means
  to Gaussian filtering.

  \part \TFT\ \ Choosing the patch to be the single pixel, $N(x) = \{x\}$, reduces non-local means
  to Yaroslavsky's filter, a special case of bilateral filtering.

  \part \TFF\ \ Increasing $h$ decreases the amount of smoothing.

  \part \TFT\ \ The averaging is restricted to a search window $W(x)$ rather than the whole image
  mainly to reduce the computational cost.
\end{parts}

\begin{solution}
  (a) FALSE --- it reduces to \emph{mean} filtering, not Gaussian filtering. If every pixel's patch
  is the whole image, all patches are identical, every weight is equal, and the weighted average
  degenerates to the plain average.

  (b) TRUE --- with a single-pixel ``patch'' the weight compares raw grey values, which is exactly
  the neighbourhood-filter family (Yaroslavsky, bilateral). This is also why NLM is more robust:
  a single noisy value is a fragile basis for similarity, a whole patch is not.

  (c) FALSE --- $h$ is the degree of filtering: larger $h$ makes the exponential weights decay more
  slowly, so more (less similar) patches contribute and the result is \emph{smoother}. $h$ should
  be tied to the noise level, which is worth estimating first.

  (d) TRUE --- and it also encodes a mild prior that useful matches are nearby. The three
  parameters of the method are exactly $h$, the patch size, and the search-window size.
\end{solution}

% =============================================================== Question 5 ===
\question
A 2D landmark-based registration uses thin-plate splines with the $N = 50$ landmarks as control
points, combined with a global affine part. How many parameters does the transformation have? How
does this compare with a naive deformable parameterization of a $256 \times 256$ image that stores
one independent displacement vector per pixel?

\begin{solution}
  Answer: the TPS model has $50 \times 2 + 6 = \mathbf{106}$ parameters; the naive field has
  $256 \times 256 \times 2 = \mathbf{131\,072}$ --- about $1\,240$ times more.

  Each control point carries a coefficient vector $d_n \in \mathbb{R}^2$ (100 parameters), and the
  2D affine part adds $4 + 2 = 6$. The comparison is the point of the question: both are
  ``deformable'' models, but the TPS parameterization builds smoothness in and needs three orders
  of magnitude fewer parameters, while the naive field is not only huge but can produce noisy,
  topology-breaking transformations for generic parameter values.
\end{solution}
\answerspace{3.6cm}

% =============================================================== Question 6 ===
\question
An Active Appearance Model extends a statistical shape model with a model of texture. Please
describe briefly (i) how the texture model is built from the training images, and (ii) how shape
and texture are combined into a single set of appearance parameters. You do not need to write down
the equations.

\begin{solution}
  Answer: \emph{(i) Building the texture model.} Every training image is \textbf{warped to the mean
  shape}, which removes the shape variation and yields a shape-free patch. The intensities are
  sampled from this shape-normalized image into a texture vector and \textbf{normalized in
  intensity} to eliminate lighting variation. PCA is then applied to the normalized texture
  vectors, exactly as it was applied to the shape coordinates: $g \approx \bar g + P_g b_g$.

  \emph{(ii) Combining shape and texture.} Each training sample is now summarized by two parameter
  vectors, $b_s$ for shape and $b_g$ for texture --- and the two are correlated (a shape change
  usually comes with a texture change). They are stacked into a single vector, with a weighting
  $W_s$ on the shape part to account for the different units of shape and texture, and a
  \textbf{second PCA} is run on the stacked vectors. The resulting parameters $c$ control shape and
  texture \emph{jointly}: a new image is generated by producing the shape-free texture from $c$ and
  warping it with the shape that the same $c$ describes.

  (Fitting then works by synthesising the model image and minimising its difference to the target,
  with the parameter-update matrix precomputed from training-set perturbations.)
\end{solution}
\answerspace{7.5cm}

% =============================================================== Question 7 ===
\question
Which of the following statements about atlas-based segmentation practice is \textbf{inaccurate}?

\begin{choices}
  \choice An atlas consists of an image together with its manual segmentation; it is registered to
  the unseen patient image and the labels are propagated along the transformation.
  \choice When an average (mean) atlas is built from a population, the choice of the reference
  subject does not influence the result --- averages built from different references are the same.
  \choice The quality of an atlas-based segmentation depends on the image similarity between atlas
  and patient, the presence of pathology, the quality of the registration, and how typical the
  patient is of the atlas population.
  \choice Shape-based averaging fuses several segmentations by averaging their distance transforms
  rather than voting on their labels.
\end{choices}

\begin{solution}
  Answer: B.

  This is the \textbf{reference-selection bias}: average models built with different reference
  subjects are demonstrably \emph{not} the same --- the lecture shows two different average images
  from two different references. The bias is precisely what motivates unbiased, group-wise
  template construction in a natural coordinate frame, rather than anchoring the average to any
  single individual.
\end{solution}

% =============================================================== Question 8 ===
\question
Which of the following statements on uncertainty estimation is \textbf{incorrect}?

\begin{choices}
  \choice One source of uncertainty is that the input $x$ may not uniquely identify the label $y$.
  \choice Another source is that the training set $D$ is itself a random sample from the
  population, and a different draw would have produced different parameters.
  \choice The Generalized Energy Distance can be evaluated on an ordinary test set with a single
  annotation per image, which is why it is the most general evaluation of uncertainty.
  \choice The Expected Calibration Error measures the average gap between the confidence a model
  assigns and the accuracy it actually achieves at that confidence.
\end{choices}

\begin{solution}
  Answer: C.

  The two evaluation routes are exactly reversed. GED (like MMD) is a \emph{distributional}
  distance between the set of model samples and the set of \textbf{observed annotations} --- it
  needs data sets with \emph{multiple labels per image}, such as multi-rater segmentations. When
  only one annotation per image is available --- the most general case --- the tool is
  \textbf{calibration}, i.e.\ ECE: whether the assigned confidence tracks the achieved accuracy.
\end{solution}

% =============================================================== Question 9 ===
\question
Which of the following statements from the machine-learning recap is \textbf{incorrect}?

\begin{choices}
  \choice Hyper-parameters --- the architecture, the stopping point of the optimization, the
  learning rate --- are chosen using a validation set, ideally as a nested optimization inside the
  training of the parameters.
  \choice The training loss is a sum over samples, so the optimal parameters do not have to satisfy
  any single sample perfectly --- they satisfy the samples on average.
  \choice The prediction error on the test set must be computed with the same differentiable loss
  that was used for training, since otherwise the numbers are not comparable.
  \choice For a new sample the prediction is $y \approx f(x; \theta^*)$, and in general
  $y \neq f(x; \theta^*)$.
\end{choices}

\begin{solution}
  Answer: C.

  Backwards on both counts. Differentiability is a requirement of \emph{training} --- the loss must
  be optimized --- not of \emph{evaluation}. On the test set one may, and often should, use a
  different, more application-appropriate function $D$, precisely in the cases where $D$ is not
  differentiable (a ranking-based measure is the lecture's example). The same logic appeared with
  segmentation: the set-based Dice score is the standard \emph{evaluation} measure even when the
  network was trained with cross-entropy or soft Dice.
\end{solution}

% ============================================================== Question 10 ===
\question
``Meaningful perturbation'' explains a classifier's decision by perturbing the input. Which of the
following statements describes the method correctly?

\begin{choices}
  \choice It searches the training set for the image that most strongly activates the class and
  presents it as the explanation.
  \choice It optimizes an information-deletion mask --- applying blur, noise or masking --- that
  covers as little of the image as possible while dropping the class score as much as possible.
  \choice It requires no evaluations of the model, which makes it the cheapest explanation method.
  \choice It computes the gradient of the class score with respect to the convolutional feature
  maps and pools it into a heatmap.
  \choice Because the mask is optimized, the resulting perturbations are guaranteed to remain
  natural-looking.
\end{choices}

\begin{solution}
  Answer: B.

  The key idea is to find the \emph{right} perturbation --- a form of information deletion --- and
  study its effect on the prediction: a mask objective balances minimal coverage against maximal
  drop in the class score. (D) describes Grad-CAM, a different method. (E) is the method's known
  caveat, inverted: unconstrained optimization can produce \emph{unnatural} perturbations ---
  adversarial-like artifacts against the explanation --- which is why the meaning of the
  explanation depends on the meaning of the changes applied to the input.
\end{solution}

% ============================================================== Question 11 ===
\question
A classification network takes grey-scale input images of size $68 \times 68$ with one input
channel. All convolutions are \emph{valid} (no padding) with stride $S = 1$:

\begin{itemize}
  \item \textbf{Conv1}: 6 kernels with $K = 5 \times 5$ \quad\textbf{Pool1}: max pooling
    $2 \times 2$, $S = 2$
  \item \textbf{Conv2}: 12 kernels with $K = 3 \times 3$ \quad\textbf{Pool2}: max pooling
    $2 \times 2$, $S = 2$
  \item \textbf{Conv3}: 24 kernels with $K = 3 \times 3$
\end{itemize}

\noindent Circle True or False:

\begin{parts}
  \part \TFF\ \ The output of Conv3 has a spatial size of $15 \times 15$.

  \part \TFT\ \ The global receptive field of a neuron in Conv3 is $20 \times 20$.

  \part \TFF\ \ The global stride of a neuron in Conv3 is $8 \times 8$.

  \part \TFT\ \ The network's convolutional layers have
  $5 \cdot 5 \cdot 1 \cdot 6 + 3 \cdot 3 \cdot 6 \cdot 12 + 3 \cdot 3 \cdot 12 \cdot 24 = 3\,390$
  kernel weights, and including one bias per kernel the total number of parameters is $3\,432$.
\end{parts}

\begin{solution}
  (a) FALSE --- $15 \times 15$ is the size \emph{entering} Conv3, not leaving it. Tracking:
  $68 \xrightarrow{\text{Conv1}} 64 \xrightarrow{\text{Pool1}} 32 \xrightarrow{\text{Conv2}} 30
  \xrightarrow{\text{Pool2}} 15 \xrightarrow{\text{Conv3}} 13$. The output is $13 \times 13$.

  (b) TRUE --- track the receptive field $R$ and jump $j$ from $R = 1$, $j = 1$ with
  $R \leftarrow R + (K-1)\cdot j$ and $j \leftarrow j \cdot S$:
  Conv1 $\rightarrow R = 5,\ j = 1$; Pool1 $\rightarrow R = 6,\ j = 2$;
  Conv2 $\rightarrow R = 10,\ j = 2$; Pool2 $\rightarrow R = 12,\ j = 4$;
  Conv3 $\rightarrow R = 20,\ j = 4$.

  (c) FALSE --- there are only \emph{two} stride-2 poolings, so the global stride is
  $2 \cdot 2 = 4$, i.e.\ $4 \times 4$. Strides compound multiplicatively over the layers that
  actually stride; Conv3 itself has stride 1 and changes nothing.

  (d) TRUE --- weights: $150 + 648 + 2\,592 = 3\,390$. Biases: one per kernel, i.e.\ per output
  channel: $6 + 12 + 24 = 42$, giving $3\,432$ in total. Pooling layers contribute neither weights
  nor biases.
\end{solution}

% ============================================================== Question 12 ===
\question
Which of the following statements are correct (multiple answers are possible)?

\begin{choices}
  \choice In FCN-style segmentation, outputs taken at several depths of the network are upsampled
  by different factors and combined, so that fine detail from shallow layers complements the
  context of deep layers.
  \choice A label map may be transformed with bilinear interpolation provided the interpolated
  values are rounded to the nearest integer label afterwards.
  \choice Matching the mean and standard deviation of two images,
  $f(j) = (j - \mu_j)\,\sigma_i/\sigma_j + \mu_i$, matches second-order statistics and is less
  sensitive to outliers than min--max normalization.
  \choice A similarity transformation has 4 degrees of freedom in 2D and 7 in 3D.
  \choice The curl of a transformation field, $\nabla \times T$, quantifies the local volumetric
  change induced by the transformation.
\end{choices}

\begin{solution}
  Answer: (A), (C) and (D) are correct.

  (B) is wrong, and the proposed repair makes it worse rather than better. Bilinear interpolation
  between labels produces values that are not labels, and rounding them does not restore meaning:
  between a voxel of label 1 and a voxel of label 3, the interpolated value $2$ rounds to label 2
  --- possibly a completely different structure that is not even adjacent to either. Label maps are
  transformed with nearest-neighbour interpolation, full stop.

  (E) attributes the wrong quantity: the \textbf{curl} carries information about the amount of
  \emph{infinitesimal rotation}. Local volumetric change is quantified by the \textbf{Jacobian
  determinant} ($\det J > 1$ expansion, $< 1$ compression, $= 1$ none), with the divergence ---
  the trace of the Jacobian --- carrying related compression/expansion information.
\end{solution}

\end{questions}

\end{document}
"
%% Compile the blank version (no answers):  pdflatex Sample_Exam_7.tex
%%
\ifx\withanswers\undefined
  \documentclass[11pt]{exam}
\else
  \documentclass[11pt,answers]{exam}
\fi

\usepackage[margin=2.4cm]{geometry}
\usepackage{amsmath,amssymb}
\usepackage{array}
\usepackage[T1]{fontenc}

% ---------------------------------------------------------------- page style
\pagestyle{foot}
\firstpageheader{}{}{}
\firstpagefooter{}{}{}
\runningheader{}{}{}
\runningfooter{}{}{Page \thepage}

% -------------------------------------------------------------- multiplechoice
% Reproduce the "(A) (B) (C) ..." labelling of the original exam.
\renewcommand{\choicelabel}{$\bigcirc$~(\Alph{choice})}

% ---------------------------------------------------------- true/false macros
% \TFT -> the statement is TRUE, \TFF -> the statement is FALSE.
% In the answer version the correct word is set in capitals and bold.
\newcommand{\TFT}{\ifprintanswers\textbf{TRUE}\hspace{1.1em}False\else True\hspace{1.1em}False\fi}
\newcommand{\TFF}{True\hspace{1.1em}\ifprintanswers\textbf{FALSE}\else False\fi}

% ------------------------------------------------------------- writing space
% Reserves room to write in the blank version; collapses in the answer version.
\newcommand{\answerspace}[1]{\ifprintanswers\else\vspace{#1}\fi}

\begin{document}

% ------------------------------------------------------------------ title page
\begin{center}
  {\LARGE\bfseries Sample Exam}\\[2mm]
  {\Large Medical Image Analysis}\\[6mm]
\end{center}

\vspace{4mm}
\noindent Student name:\hrulefill

\vspace{6mm}
\noindent Immatriculation number:\hrulefill

\vspace{10mm}

\begin{center}
  {\large\bfseries Exam Questions}
\end{center}

\noindent\textbf{Note:} All the questions have equal weight.

\vspace{4mm}

\begin{questions}

% =============================================================== Question 1 ===
\question
In 2D acquisitions, slices are acquired in one pre-defined plane and stacked along the remaining
direction: axial slices stack head-to-toe, sagittal slices left-to-right, coronal slices
anterior-to-posterior. Voxel sizes are given as in-plane (two values) and through-plane (the slice
thickness), and the field-of-view is image size $\times$ voxel size.

\vspace{1mm}
\noindent You are handed a volume with the following properties, but the acquisition plane is not
recorded: field-of-view $240$~mm head-to-toe, $240$~mm left-to-right, $180$~mm
anterior-to-posterior; voxel size in-plane $0.75 \times 0.75$~mm$^2$, through-plane $4.5$~mm.

\begin{parts}

  \part In which plane was this volume acquired? Justify your answer from the numbers.

  \begin{solution}
    Answer: a \textbf{coronal} acquisition.

    The through-plane direction is the one whose extent is sampled at $4.5$~mm, and only
    anterior-to-posterior works: $180 / 4.5 = 40$ slices, an integer. If either of the other two
    directions were through-plane, it would need $240 / 4.5 = 53.3$ slices --- not an integer, so
    the numbers are inconsistent with any other assignment. Slices stacked along
    anterior-to-posterior lie in the coronal plane.
  \end{solution}
\answerspace{3.0cm}

  \part How many voxels does the volume have in the head-to-toe, left-to-right and
  anterior-to-posterior directions, respectively?

  \begin{solution}
    Answer: $320 \times 320 \times 40$~voxels.

    Head-to-toe and left-to-right are in-plane: $240 / 0.75 = 320$ each. Anterior-to-posterior is
    through-plane: $180 / 4.5 = 40$.
  \end{solution}
\answerspace{2.2cm}

  \part The volume is resampled to isotropic $1.5 \times 1.5 \times 1.5$~mm$^3$ voxels. Give the
  resize factors and the resulting image size, and verify that the field-of-view is unchanged.

  \begin{solution}
    Answer: resize factors $0.5 \times 0.5 \times 3.0$, resulting image size
    $160 \times 160 \times 120$~voxels.

    The factor is the ratio of old to new voxel size: $0.75/1.5 = 0.5$ in-plane,
    $4.5/1.5 = 3.0$ through-plane. Sizes: $320 \cdot 0.5 = 160$, $320 \cdot 0.5 = 160$,
    $40 \cdot 3.0 = 120$. Field-of-view check: $160 \cdot 1.5 = 240$, $160 \cdot 1.5 = 240$,
    $120 \cdot 1.5 = 180$~mm --- unchanged, as it must be. (The anterior-to-posterior direction is
    upsampled by 3: two of every three samples there are interpolated, not measured.)
  \end{solution}
\answerspace{3.6cm}

\end{parts}

% =============================================================== Question 2 ===
\question
In bilinear interpolation the intensity at a query point $x_1$ (with $x_1^0 \le x_1 \le x_1^1$) is
first interpolated along one dimension from the two neighbouring pixels. Which of the two
expressions below is the correct linear interpolation? Please circle the correct choice.

\vspace{2mm}
\begin{center}
  $\displaystyle I(x_1) = \frac{x_1^1 - x_1}{x_1^1 - x_1^0}\, I(x_1^0)
  + \frac{x_1 - x_1^0}{x_1^1 - x_1^0}\, I(x_1^1)$
  \hspace{1.1cm}
  $\displaystyle I(x_1) = \frac{x_1 - x_1^0}{x_1^1 - x_1^0}\, I(x_1^0)
  + \frac{x_1^1 - x_1}{x_1^1 - x_1^0}\, I(x_1^1)$
\end{center}
\vspace{2mm}

\begin{solution}
  Answer: Left one.

  Each pixel's weight must be proportional to the query point's distance from the \emph{other}
  pixel --- nearness earns weight. The sanity check takes two seconds: set $x_1 = x_1^0$. The left
  expression gives $1 \cdot I(x_1^0) + 0 \cdot I(x_1^1) = I(x_1^0)$, as an interpolant must. The
  right expression gives $I(x_1^1)$ --- the wrong pixel entirely; its weights are inverted.
  Bilinear interpolation applies this step along $x_1$ and then along $x_2$ (trilinear adds the
  third dimension, using the 8 cube corners).
\end{solution}
\answerspace{1.5cm}

% =============================================================== Question 3 ===
\question
Probabilistic reasoning underlies most of this course.

\begin{parts}
  \part If each pixel intensity is treated as a realization of a random variable, the (normalized)
  histogram and the cumulative histogram of an image approximate two standard probabilistic
  objects. Which two --- and what assumption about the pixels makes this approximation valid?

  \begin{solution}
    Answer: the normalized histogram approximates the \textbf{probability density function} of the
    intensity, and the cumulative histogram approximates the \textbf{cumulative distribution
    function}.

    The assumption is that each pixel intensity is an \textbf{independent realization} of the same
    random variable $I$ --- only then is counting pixels per intensity bin an estimate of the
    underlying distribution.
  \end{solution}
\answerspace{2.8cm}

  \part The Maximum-A-Posteriori estimate is $\arg\max_l\, p(i \mid l)\, p(l)$ and the Maximum
  Likelihood estimate is $\arg\max_l\, p(i \mid l)$. Under what condition on the prior are the two
  identical?

  \begin{solution}
    Answer: when the prior is \textbf{uniform}, $p(l) = c$ for all $l$.

    A constant factor does not change the location of a maximum, so the prior drops out of the
    $\arg\max$ and MAP reduces to MLE. Conversely, whenever anyone uses MLE they are implicitly
    assuming that no value of $l$ is more plausible than any other before seeing the data.
  \end{solution}
\answerspace{2.2cm}
\end{parts}

% =============================================================== Question 4 ===
\question
Non-local means computes $\hat I(x) = \sum_{y} w(x,y) J(y)$ with weights from the similarity of
the patches $N(x)$ and $N(y)$, a degree-of-filtering parameter $h$, and averaging restricted to a
search window $W(x)$. Please indicate whether the statement is true or false by circling the
correct choice in each of the following.

\begin{parts}
  \part \TFF\ \ Choosing the patch to be the entire image, $N(x) = \Omega$, reduces non-local means
  to Gaussian filtering.

  \part \TFT\ \ Choosing the patch to be the single pixel, $N(x) = \{x\}$, reduces non-local means
  to Yaroslavsky's filter, a special case of bilateral filtering.

  \part \TFF\ \ Increasing $h$ decreases the amount of smoothing.

  \part \TFT\ \ The averaging is restricted to a search window $W(x)$ rather than the whole image
  mainly to reduce the computational cost.
\end{parts}

\begin{solution}
  (a) FALSE --- it reduces to \emph{mean} filtering, not Gaussian filtering. If every pixel's patch
  is the whole image, all patches are identical, every weight is equal, and the weighted average
  degenerates to the plain average.

  (b) TRUE --- with a single-pixel ``patch'' the weight compares raw grey values, which is exactly
  the neighbourhood-filter family (Yaroslavsky, bilateral). This is also why NLM is more robust:
  a single noisy value is a fragile basis for similarity, a whole patch is not.

  (c) FALSE --- $h$ is the degree of filtering: larger $h$ makes the exponential weights decay more
  slowly, so more (less similar) patches contribute and the result is \emph{smoother}. $h$ should
  be tied to the noise level, which is worth estimating first.

  (d) TRUE --- and it also encodes a mild prior that useful matches are nearby. The three
  parameters of the method are exactly $h$, the patch size, and the search-window size.
\end{solution}

% =============================================================== Question 5 ===
\question
A 2D landmark-based registration uses thin-plate splines with the $N = 50$ landmarks as control
points, combined with a global affine part. How many parameters does the transformation have? How
does this compare with a naive deformable parameterization of a $256 \times 256$ image that stores
one independent displacement vector per pixel?

\begin{solution}
  Answer: the TPS model has $50 \times 2 + 6 = \mathbf{106}$ parameters; the naive field has
  $256 \times 256 \times 2 = \mathbf{131\,072}$ --- about $1\,240$ times more.

  Each control point carries a coefficient vector $d_n \in \mathbb{R}^2$ (100 parameters), and the
  2D affine part adds $4 + 2 = 6$. The comparison is the point of the question: both are
  ``deformable'' models, but the TPS parameterization builds smoothness in and needs three orders
  of magnitude fewer parameters, while the naive field is not only huge but can produce noisy,
  topology-breaking transformations for generic parameter values.
\end{solution}
\answerspace{3.6cm}

% =============================================================== Question 6 ===
\question
An Active Appearance Model extends a statistical shape model with a model of texture. Please
describe briefly (i) how the texture model is built from the training images, and (ii) how shape
and texture are combined into a single set of appearance parameters. You do not need to write down
the equations.

\begin{solution}
  Answer: \emph{(i) Building the texture model.} Every training image is \textbf{warped to the mean
  shape}, which removes the shape variation and yields a shape-free patch. The intensities are
  sampled from this shape-normalized image into a texture vector and \textbf{normalized in
  intensity} to eliminate lighting variation. PCA is then applied to the normalized texture
  vectors, exactly as it was applied to the shape coordinates: $g \approx \bar g + P_g b_g$.

  \emph{(ii) Combining shape and texture.} Each training sample is now summarized by two parameter
  vectors, $b_s$ for shape and $b_g$ for texture --- and the two are correlated (a shape change
  usually comes with a texture change). They are stacked into a single vector, with a weighting
  $W_s$ on the shape part to account for the different units of shape and texture, and a
  \textbf{second PCA} is run on the stacked vectors. The resulting parameters $c$ control shape and
  texture \emph{jointly}: a new image is generated by producing the shape-free texture from $c$ and
  warping it with the shape that the same $c$ describes.

  (Fitting then works by synthesising the model image and minimising its difference to the target,
  with the parameter-update matrix precomputed from training-set perturbations.)
\end{solution}
\answerspace{7.5cm}

% =============================================================== Question 7 ===
\question
Which of the following statements about atlas-based segmentation practice is \textbf{inaccurate}?

\begin{choices}
  \choice An atlas consists of an image together with its manual segmentation; it is registered to
  the unseen patient image and the labels are propagated along the transformation.
  \choice When an average (mean) atlas is built from a population, the choice of the reference
  subject does not influence the result --- averages built from different references are the same.
  \choice The quality of an atlas-based segmentation depends on the image similarity between atlas
  and patient, the presence of pathology, the quality of the registration, and how typical the
  patient is of the atlas population.
  \choice Shape-based averaging fuses several segmentations by averaging their distance transforms
  rather than voting on their labels.
\end{choices}

\begin{solution}
  Answer: B.

  This is the \textbf{reference-selection bias}: average models built with different reference
  subjects are demonstrably \emph{not} the same --- the lecture shows two different average images
  from two different references. The bias is precisely what motivates unbiased, group-wise
  template construction in a natural coordinate frame, rather than anchoring the average to any
  single individual.
\end{solution}

% =============================================================== Question 8 ===
\question
Which of the following statements on uncertainty estimation is \textbf{incorrect}?

\begin{choices}
  \choice One source of uncertainty is that the input $x$ may not uniquely identify the label $y$.
  \choice Another source is that the training set $D$ is itself a random sample from the
  population, and a different draw would have produced different parameters.
  \choice The Generalized Energy Distance can be evaluated on an ordinary test set with a single
  annotation per image, which is why it is the most general evaluation of uncertainty.
  \choice The Expected Calibration Error measures the average gap between the confidence a model
  assigns and the accuracy it actually achieves at that confidence.
\end{choices}

\begin{solution}
  Answer: C.

  The two evaluation routes are exactly reversed. GED (like MMD) is a \emph{distributional}
  distance between the set of model samples and the set of \textbf{observed annotations} --- it
  needs data sets with \emph{multiple labels per image}, such as multi-rater segmentations. When
  only one annotation per image is available --- the most general case --- the tool is
  \textbf{calibration}, i.e.\ ECE: whether the assigned confidence tracks the achieved accuracy.
\end{solution}

% =============================================================== Question 9 ===
\question
Which of the following statements from the machine-learning recap is \textbf{incorrect}?

\begin{choices}
  \choice Hyper-parameters --- the architecture, the stopping point of the optimization, the
  learning rate --- are chosen using a validation set, ideally as a nested optimization inside the
  training of the parameters.
  \choice The training loss is a sum over samples, so the optimal parameters do not have to satisfy
  any single sample perfectly --- they satisfy the samples on average.
  \choice The prediction error on the test set must be computed with the same differentiable loss
  that was used for training, since otherwise the numbers are not comparable.
  \choice For a new sample the prediction is $y \approx f(x; \theta^*)$, and in general
  $y \neq f(x; \theta^*)$.
\end{choices}

\begin{solution}
  Answer: C.

  Backwards on both counts. Differentiability is a requirement of \emph{training} --- the loss must
  be optimized --- not of \emph{evaluation}. On the test set one may, and often should, use a
  different, more application-appropriate function $D$, precisely in the cases where $D$ is not
  differentiable (a ranking-based measure is the lecture's example). The same logic appeared with
  segmentation: the set-based Dice score is the standard \emph{evaluation} measure even when the
  network was trained with cross-entropy or soft Dice.
\end{solution}

% ============================================================== Question 10 ===
\question
``Meaningful perturbation'' explains a classifier's decision by perturbing the input. Which of the
following statements describes the method correctly?

\begin{choices}
  \choice It searches the training set for the image that most strongly activates the class and
  presents it as the explanation.
  \choice It optimizes an information-deletion mask --- applying blur, noise or masking --- that
  covers as little of the image as possible while dropping the class score as much as possible.
  \choice It requires no evaluations of the model, which makes it the cheapest explanation method.
  \choice It computes the gradient of the class score with respect to the convolutional feature
  maps and pools it into a heatmap.
  \choice Because the mask is optimized, the resulting perturbations are guaranteed to remain
  natural-looking.
\end{choices}

\begin{solution}
  Answer: B.

  The key idea is to find the \emph{right} perturbation --- a form of information deletion --- and
  study its effect on the prediction: a mask objective balances minimal coverage against maximal
  drop in the class score. (D) describes Grad-CAM, a different method. (E) is the method's known
  caveat, inverted: unconstrained optimization can produce \emph{unnatural} perturbations ---
  adversarial-like artifacts against the explanation --- which is why the meaning of the
  explanation depends on the meaning of the changes applied to the input.
\end{solution}

% ============================================================== Question 11 ===
\question
A classification network takes grey-scale input images of size $68 \times 68$ with one input
channel. All convolutions are \emph{valid} (no padding) with stride $S = 1$:

\begin{itemize}
  \item \textbf{Conv1}: 6 kernels with $K = 5 \times 5$ \quad\textbf{Pool1}: max pooling
    $2 \times 2$, $S = 2$
  \item \textbf{Conv2}: 12 kernels with $K = 3 \times 3$ \quad\textbf{Pool2}: max pooling
    $2 \times 2$, $S = 2$
  \item \textbf{Conv3}: 24 kernels with $K = 3 \times 3$
\end{itemize}

\noindent Circle True or False:

\begin{parts}
  \part \TFF\ \ The output of Conv3 has a spatial size of $15 \times 15$.

  \part \TFT\ \ The global receptive field of a neuron in Conv3 is $20 \times 20$.

  \part \TFF\ \ The global stride of a neuron in Conv3 is $8 \times 8$.

  \part \TFT\ \ The network's convolutional layers have
  $5 \cdot 5 \cdot 1 \cdot 6 + 3 \cdot 3 \cdot 6 \cdot 12 + 3 \cdot 3 \cdot 12 \cdot 24 = 3\,390$
  kernel weights, and including one bias per kernel the total number of parameters is $3\,432$.
\end{parts}

\begin{solution}
  (a) FALSE --- $15 \times 15$ is the size \emph{entering} Conv3, not leaving it. Tracking:
  $68 \xrightarrow{\text{Conv1}} 64 \xrightarrow{\text{Pool1}} 32 \xrightarrow{\text{Conv2}} 30
  \xrightarrow{\text{Pool2}} 15 \xrightarrow{\text{Conv3}} 13$. The output is $13 \times 13$.

  (b) TRUE --- track the receptive field $R$ and jump $j$ from $R = 1$, $j = 1$ with
  $R \leftarrow R + (K-1)\cdot j$ and $j \leftarrow j \cdot S$:
  Conv1 $\rightarrow R = 5,\ j = 1$; Pool1 $\rightarrow R = 6,\ j = 2$;
  Conv2 $\rightarrow R = 10,\ j = 2$; Pool2 $\rightarrow R = 12,\ j = 4$;
  Conv3 $\rightarrow R = 20,\ j = 4$.

  (c) FALSE --- there are only \emph{two} stride-2 poolings, so the global stride is
  $2 \cdot 2 = 4$, i.e.\ $4 \times 4$. Strides compound multiplicatively over the layers that
  actually stride; Conv3 itself has stride 1 and changes nothing.

  (d) TRUE --- weights: $150 + 648 + 2\,592 = 3\,390$. Biases: one per kernel, i.e.\ per output
  channel: $6 + 12 + 24 = 42$, giving $3\,432$ in total. Pooling layers contribute neither weights
  nor biases.
\end{solution}

% ============================================================== Question 12 ===
\question
Which of the following statements are correct (multiple answers are possible)?

\begin{choices}
  \choice In FCN-style segmentation, outputs taken at several depths of the network are upsampled
  by different factors and combined, so that fine detail from shallow layers complements the
  context of deep layers.
  \choice A label map may be transformed with bilinear interpolation provided the interpolated
  values are rounded to the nearest integer label afterwards.
  \choice Matching the mean and standard deviation of two images,
  $f(j) = (j - \mu_j)\,\sigma_i/\sigma_j + \mu_i$, matches second-order statistics and is less
  sensitive to outliers than min--max normalization.
  \choice A similarity transformation has 4 degrees of freedom in 2D and 7 in 3D.
  \choice The curl of a transformation field, $\nabla \times T$, quantifies the local volumetric
  change induced by the transformation.
\end{choices}

\begin{solution}
  Answer: (A), (C) and (D) are correct.

  (B) is wrong, and the proposed repair makes it worse rather than better. Bilinear interpolation
  between labels produces values that are not labels, and rounding them does not restore meaning:
  between a voxel of label 1 and a voxel of label 3, the interpolated value $2$ rounds to label 2
  --- possibly a completely different structure that is not even adjacent to either. Label maps are
  transformed with nearest-neighbour interpolation, full stop.

  (E) attributes the wrong quantity: the \textbf{curl} carries information about the amount of
  \emph{infinitesimal rotation}. Local volumetric change is quantified by the \textbf{Jacobian
  determinant} ($\det J > 1$ expansion, $< 1$ compression, $= 1$ none), with the divergence ---
  the trace of the Jacobian --- carrying related compression/expansion information.
\end{solution}

\end{questions}

\end{document}
"
%% Compile the blank version (no answers):  pdflatex Sample_Exam_7.tex
%%
\ifx\withanswers\undefined
  \documentclass[11pt]{exam}
\else
  \documentclass[11pt,answers]{exam}
\fi

\usepackage[margin=2.4cm]{geometry}
\usepackage{amsmath,amssymb}
\usepackage{array}
\usepackage[T1]{fontenc}

% ---------------------------------------------------------------- page style
\pagestyle{foot}
\firstpageheader{}{}{}
\firstpagefooter{}{}{}
\runningheader{}{}{}
\runningfooter{}{}{Page \thepage}

% -------------------------------------------------------------- multiplechoice
% Reproduce the "(A) (B) (C) ..." labelling of the original exam.
\renewcommand{\choicelabel}{$\bigcirc$~(\Alph{choice})}

% ---------------------------------------------------------- true/false macros
% \TFT -> the statement is TRUE, \TFF -> the statement is FALSE.
% In the answer version the correct word is set in capitals and bold.
\newcommand{\TFT}{\ifprintanswers\textbf{TRUE}\hspace{1.1em}False\else True\hspace{1.1em}False\fi}
\newcommand{\TFF}{True\hspace{1.1em}\ifprintanswers\textbf{FALSE}\else False\fi}

% ------------------------------------------------------------- writing space
% Reserves room to write in the blank version; collapses in the answer version.
\newcommand{\answerspace}[1]{\ifprintanswers\else\vspace{#1}\fi}

\begin{document}

% ------------------------------------------------------------------ title page
\begin{center}
  {\LARGE\bfseries Sample Exam}\\[2mm]
  {\Large Medical Image Analysis}\\[6mm]
\end{center}

\vspace{4mm}
\noindent Student name:\hrulefill

\vspace{6mm}
\noindent Immatriculation number:\hrulefill

\vspace{10mm}

\begin{center}
  {\large\bfseries Exam Questions}
\end{center}

\noindent\textbf{Note:} All the questions have equal weight.

\vspace{4mm}

\begin{questions}

% =============================================================== Question 1 ===
\question
In 2D acquisitions, slices are acquired in one pre-defined plane and stacked along the remaining
direction: axial slices stack head-to-toe, sagittal slices left-to-right, coronal slices
anterior-to-posterior. Voxel sizes are given as in-plane (two values) and through-plane (the slice
thickness), and the field-of-view is image size $\times$ voxel size.

\vspace{1mm}
\noindent You are handed a volume with the following properties, but the acquisition plane is not
recorded: field-of-view $240$~mm head-to-toe, $240$~mm left-to-right, $180$~mm
anterior-to-posterior; voxel size in-plane $0.75 \times 0.75$~mm$^2$, through-plane $4.5$~mm.

\begin{parts}

  \part In which plane was this volume acquired? Justify your answer from the numbers.

  \begin{solution}
    Answer: a \textbf{coronal} acquisition.

    The through-plane direction is the one whose extent is sampled at $4.5$~mm, and only
    anterior-to-posterior works: $180 / 4.5 = 40$ slices, an integer. If either of the other two
    directions were through-plane, it would need $240 / 4.5 = 53.3$ slices --- not an integer, so
    the numbers are inconsistent with any other assignment. Slices stacked along
    anterior-to-posterior lie in the coronal plane.
  \end{solution}
\answerspace{3.0cm}

  \part How many voxels does the volume have in the head-to-toe, left-to-right and
  anterior-to-posterior directions, respectively?

  \begin{solution}
    Answer: $320 \times 320 \times 40$~voxels.

    Head-to-toe and left-to-right are in-plane: $240 / 0.75 = 320$ each. Anterior-to-posterior is
    through-plane: $180 / 4.5 = 40$.
  \end{solution}
\answerspace{2.2cm}

  \part The volume is resampled to isotropic $1.5 \times 1.5 \times 1.5$~mm$^3$ voxels. Give the
  resize factors and the resulting image size, and verify that the field-of-view is unchanged.

  \begin{solution}
    Answer: resize factors $0.5 \times 0.5 \times 3.0$, resulting image size
    $160 \times 160 \times 120$~voxels.

    The factor is the ratio of old to new voxel size: $0.75/1.5 = 0.5$ in-plane,
    $4.5/1.5 = 3.0$ through-plane. Sizes: $320 \cdot 0.5 = 160$, $320 \cdot 0.5 = 160$,
    $40 \cdot 3.0 = 120$. Field-of-view check: $160 \cdot 1.5 = 240$, $160 \cdot 1.5 = 240$,
    $120 \cdot 1.5 = 180$~mm --- unchanged, as it must be. (The anterior-to-posterior direction is
    upsampled by 3: two of every three samples there are interpolated, not measured.)
  \end{solution}
\answerspace{3.6cm}

\end{parts}

% =============================================================== Question 2 ===
\question
In bilinear interpolation the intensity at a query point $x_1$ (with $x_1^0 \le x_1 \le x_1^1$) is
first interpolated along one dimension from the two neighbouring pixels. Which of the two
expressions below is the correct linear interpolation? Please circle the correct choice.

\vspace{2mm}
\begin{center}
  $\displaystyle I(x_1) = \frac{x_1^1 - x_1}{x_1^1 - x_1^0}\, I(x_1^0)
  + \frac{x_1 - x_1^0}{x_1^1 - x_1^0}\, I(x_1^1)$
  \hspace{1.1cm}
  $\displaystyle I(x_1) = \frac{x_1 - x_1^0}{x_1^1 - x_1^0}\, I(x_1^0)
  + \frac{x_1^1 - x_1}{x_1^1 - x_1^0}\, I(x_1^1)$
\end{center}
\vspace{2mm}

\begin{solution}
  Answer: Left one.

  Each pixel's weight must be proportional to the query point's distance from the \emph{other}
  pixel --- nearness earns weight. The sanity check takes two seconds: set $x_1 = x_1^0$. The left
  expression gives $1 \cdot I(x_1^0) + 0 \cdot I(x_1^1) = I(x_1^0)$, as an interpolant must. The
  right expression gives $I(x_1^1)$ --- the wrong pixel entirely; its weights are inverted.
  Bilinear interpolation applies this step along $x_1$ and then along $x_2$ (trilinear adds the
  third dimension, using the 8 cube corners).
\end{solution}
\answerspace{1.5cm}

% =============================================================== Question 3 ===
\question
Probabilistic reasoning underlies most of this course.

\begin{parts}
  \part If each pixel intensity is treated as a realization of a random variable, the (normalized)
  histogram and the cumulative histogram of an image approximate two standard probabilistic
  objects. Which two --- and what assumption about the pixels makes this approximation valid?

  \begin{solution}
    Answer: the normalized histogram approximates the \textbf{probability density function} of the
    intensity, and the cumulative histogram approximates the \textbf{cumulative distribution
    function}.

    The assumption is that each pixel intensity is an \textbf{independent realization} of the same
    random variable $I$ --- only then is counting pixels per intensity bin an estimate of the
    underlying distribution.
  \end{solution}
\answerspace{2.8cm}

  \part The Maximum-A-Posteriori estimate is $\arg\max_l\, p(i \mid l)\, p(l)$ and the Maximum
  Likelihood estimate is $\arg\max_l\, p(i \mid l)$. Under what condition on the prior are the two
  identical?

  \begin{solution}
    Answer: when the prior is \textbf{uniform}, $p(l) = c$ for all $l$.

    A constant factor does not change the location of a maximum, so the prior drops out of the
    $\arg\max$ and MAP reduces to MLE. Conversely, whenever anyone uses MLE they are implicitly
    assuming that no value of $l$ is more plausible than any other before seeing the data.
  \end{solution}
\answerspace{2.2cm}
\end{parts}

% =============================================================== Question 4 ===
\question
Non-local means computes $\hat I(x) = \sum_{y} w(x,y) J(y)$ with weights from the similarity of
the patches $N(x)$ and $N(y)$, a degree-of-filtering parameter $h$, and averaging restricted to a
search window $W(x)$. Please indicate whether the statement is true or false by circling the
correct choice in each of the following.

\begin{parts}
  \part \TFF\ \ Choosing the patch to be the entire image, $N(x) = \Omega$, reduces non-local means
  to Gaussian filtering.

  \part \TFT\ \ Choosing the patch to be the single pixel, $N(x) = \{x\}$, reduces non-local means
  to Yaroslavsky's filter, a special case of bilateral filtering.

  \part \TFF\ \ Increasing $h$ decreases the amount of smoothing.

  \part \TFT\ \ The averaging is restricted to a search window $W(x)$ rather than the whole image
  mainly to reduce the computational cost.
\end{parts}

\begin{solution}
  (a) FALSE --- it reduces to \emph{mean} filtering, not Gaussian filtering. If every pixel's patch
  is the whole image, all patches are identical, every weight is equal, and the weighted average
  degenerates to the plain average.

  (b) TRUE --- with a single-pixel ``patch'' the weight compares raw grey values, which is exactly
  the neighbourhood-filter family (Yaroslavsky, bilateral). This is also why NLM is more robust:
  a single noisy value is a fragile basis for similarity, a whole patch is not.

  (c) FALSE --- $h$ is the degree of filtering: larger $h$ makes the exponential weights decay more
  slowly, so more (less similar) patches contribute and the result is \emph{smoother}. $h$ should
  be tied to the noise level, which is worth estimating first.

  (d) TRUE --- and it also encodes a mild prior that useful matches are nearby. The three
  parameters of the method are exactly $h$, the patch size, and the search-window size.
\end{solution}

% =============================================================== Question 5 ===
\question
A 2D landmark-based registration uses thin-plate splines with the $N = 50$ landmarks as control
points, combined with a global affine part. How many parameters does the transformation have? How
does this compare with a naive deformable parameterization of a $256 \times 256$ image that stores
one independent displacement vector per pixel?

\begin{solution}
  Answer: the TPS model has $50 \times 2 + 6 = \mathbf{106}$ parameters; the naive field has
  $256 \times 256 \times 2 = \mathbf{131\,072}$ --- about $1\,240$ times more.

  Each control point carries a coefficient vector $d_n \in \mathbb{R}^2$ (100 parameters), and the
  2D affine part adds $4 + 2 = 6$. The comparison is the point of the question: both are
  ``deformable'' models, but the TPS parameterization builds smoothness in and needs three orders
  of magnitude fewer parameters, while the naive field is not only huge but can produce noisy,
  topology-breaking transformations for generic parameter values.
\end{solution}
\answerspace{3.6cm}

% =============================================================== Question 6 ===
\question
An Active Appearance Model extends a statistical shape model with a model of texture. Please
describe briefly (i) how the texture model is built from the training images, and (ii) how shape
and texture are combined into a single set of appearance parameters. You do not need to write down
the equations.

\begin{solution}
  Answer: \emph{(i) Building the texture model.} Every training image is \textbf{warped to the mean
  shape}, which removes the shape variation and yields a shape-free patch. The intensities are
  sampled from this shape-normalized image into a texture vector and \textbf{normalized in
  intensity} to eliminate lighting variation. PCA is then applied to the normalized texture
  vectors, exactly as it was applied to the shape coordinates: $g \approx \bar g + P_g b_g$.

  \emph{(ii) Combining shape and texture.} Each training sample is now summarized by two parameter
  vectors, $b_s$ for shape and $b_g$ for texture --- and the two are correlated (a shape change
  usually comes with a texture change). They are stacked into a single vector, with a weighting
  $W_s$ on the shape part to account for the different units of shape and texture, and a
  \textbf{second PCA} is run on the stacked vectors. The resulting parameters $c$ control shape and
  texture \emph{jointly}: a new image is generated by producing the shape-free texture from $c$ and
  warping it with the shape that the same $c$ describes.

  (Fitting then works by synthesising the model image and minimising its difference to the target,
  with the parameter-update matrix precomputed from training-set perturbations.)
\end{solution}
\answerspace{7.5cm}

% =============================================================== Question 7 ===
\question
Which of the following statements about atlas-based segmentation practice is \textbf{inaccurate}?

\begin{choices}
  \choice An atlas consists of an image together with its manual segmentation; it is registered to
  the unseen patient image and the labels are propagated along the transformation.
  \choice When an average (mean) atlas is built from a population, the choice of the reference
  subject does not influence the result --- averages built from different references are the same.
  \choice The quality of an atlas-based segmentation depends on the image similarity between atlas
  and patient, the presence of pathology, the quality of the registration, and how typical the
  patient is of the atlas population.
  \choice Shape-based averaging fuses several segmentations by averaging their distance transforms
  rather than voting on their labels.
\end{choices}

\begin{solution}
  Answer: B.

  This is the \textbf{reference-selection bias}: average models built with different reference
  subjects are demonstrably \emph{not} the same --- the lecture shows two different average images
  from two different references. The bias is precisely what motivates unbiased, group-wise
  template construction in a natural coordinate frame, rather than anchoring the average to any
  single individual.
\end{solution}

% =============================================================== Question 8 ===
\question
Which of the following statements on uncertainty estimation is \textbf{incorrect}?

\begin{choices}
  \choice One source of uncertainty is that the input $x$ may not uniquely identify the label $y$.
  \choice Another source is that the training set $D$ is itself a random sample from the
  population, and a different draw would have produced different parameters.
  \choice The Generalized Energy Distance can be evaluated on an ordinary test set with a single
  annotation per image, which is why it is the most general evaluation of uncertainty.
  \choice The Expected Calibration Error measures the average gap between the confidence a model
  assigns and the accuracy it actually achieves at that confidence.
\end{choices}

\begin{solution}
  Answer: C.

  The two evaluation routes are exactly reversed. GED (like MMD) is a \emph{distributional}
  distance between the set of model samples and the set of \textbf{observed annotations} --- it
  needs data sets with \emph{multiple labels per image}, such as multi-rater segmentations. When
  only one annotation per image is available --- the most general case --- the tool is
  \textbf{calibration}, i.e.\ ECE: whether the assigned confidence tracks the achieved accuracy.
\end{solution}

% =============================================================== Question 9 ===
\question
Which of the following statements from the machine-learning recap is \textbf{incorrect}?

\begin{choices}
  \choice Hyper-parameters --- the architecture, the stopping point of the optimization, the
  learning rate --- are chosen using a validation set, ideally as a nested optimization inside the
  training of the parameters.
  \choice The training loss is a sum over samples, so the optimal parameters do not have to satisfy
  any single sample perfectly --- they satisfy the samples on average.
  \choice The prediction error on the test set must be computed with the same differentiable loss
  that was used for training, since otherwise the numbers are not comparable.
  \choice For a new sample the prediction is $y \approx f(x; \theta^*)$, and in general
  $y \neq f(x; \theta^*)$.
\end{choices}

\begin{solution}
  Answer: C.

  Backwards on both counts. Differentiability is a requirement of \emph{training} --- the loss must
  be optimized --- not of \emph{evaluation}. On the test set one may, and often should, use a
  different, more application-appropriate function $D$, precisely in the cases where $D$ is not
  differentiable (a ranking-based measure is the lecture's example). The same logic appeared with
  segmentation: the set-based Dice score is the standard \emph{evaluation} measure even when the
  network was trained with cross-entropy or soft Dice.
\end{solution}

% ============================================================== Question 10 ===
\question
``Meaningful perturbation'' explains a classifier's decision by perturbing the input. Which of the
following statements describes the method correctly?

\begin{choices}
  \choice It searches the training set for the image that most strongly activates the class and
  presents it as the explanation.
  \choice It optimizes an information-deletion mask --- applying blur, noise or masking --- that
  covers as little of the image as possible while dropping the class score as much as possible.
  \choice It requires no evaluations of the model, which makes it the cheapest explanation method.
  \choice It computes the gradient of the class score with respect to the convolutional feature
  maps and pools it into a heatmap.
  \choice Because the mask is optimized, the resulting perturbations are guaranteed to remain
  natural-looking.
\end{choices}

\begin{solution}
  Answer: B.

  The key idea is to find the \emph{right} perturbation --- a form of information deletion --- and
  study its effect on the prediction: a mask objective balances minimal coverage against maximal
  drop in the class score. (D) describes Grad-CAM, a different method. (E) is the method's known
  caveat, inverted: unconstrained optimization can produce \emph{unnatural} perturbations ---
  adversarial-like artifacts against the explanation --- which is why the meaning of the
  explanation depends on the meaning of the changes applied to the input.
\end{solution}

% ============================================================== Question 11 ===
\question
A classification network takes grey-scale input images of size $68 \times 68$ with one input
channel. All convolutions are \emph{valid} (no padding) with stride $S = 1$:

\begin{itemize}
  \item \textbf{Conv1}: 6 kernels with $K = 5 \times 5$ \quad\textbf{Pool1}: max pooling
    $2 \times 2$, $S = 2$
  \item \textbf{Conv2}: 12 kernels with $K = 3 \times 3$ \quad\textbf{Pool2}: max pooling
    $2 \times 2$, $S = 2$
  \item \textbf{Conv3}: 24 kernels with $K = 3 \times 3$
\end{itemize}

\noindent Circle True or False:

\begin{parts}
  \part \TFF\ \ The output of Conv3 has a spatial size of $15 \times 15$.

  \part \TFT\ \ The global receptive field of a neuron in Conv3 is $20 \times 20$.

  \part \TFF\ \ The global stride of a neuron in Conv3 is $8 \times 8$.

  \part \TFT\ \ The network's convolutional layers have
  $5 \cdot 5 \cdot 1 \cdot 6 + 3 \cdot 3 \cdot 6 \cdot 12 + 3 \cdot 3 \cdot 12 \cdot 24 = 3\,390$
  kernel weights, and including one bias per kernel the total number of parameters is $3\,432$.
\end{parts}

\begin{solution}
  (a) FALSE --- $15 \times 15$ is the size \emph{entering} Conv3, not leaving it. Tracking:
  $68 \xrightarrow{\text{Conv1}} 64 \xrightarrow{\text{Pool1}} 32 \xrightarrow{\text{Conv2}} 30
  \xrightarrow{\text{Pool2}} 15 \xrightarrow{\text{Conv3}} 13$. The output is $13 \times 13$.

  (b) TRUE --- track the receptive field $R$ and jump $j$ from $R = 1$, $j = 1$ with
  $R \leftarrow R + (K-1)\cdot j$ and $j \leftarrow j \cdot S$:
  Conv1 $\rightarrow R = 5,\ j = 1$; Pool1 $\rightarrow R = 6,\ j = 2$;
  Conv2 $\rightarrow R = 10,\ j = 2$; Pool2 $\rightarrow R = 12,\ j = 4$;
  Conv3 $\rightarrow R = 20,\ j = 4$.

  (c) FALSE --- there are only \emph{two} stride-2 poolings, so the global stride is
  $2 \cdot 2 = 4$, i.e.\ $4 \times 4$. Strides compound multiplicatively over the layers that
  actually stride; Conv3 itself has stride 1 and changes nothing.

  (d) TRUE --- weights: $150 + 648 + 2\,592 = 3\,390$. Biases: one per kernel, i.e.\ per output
  channel: $6 + 12 + 24 = 42$, giving $3\,432$ in total. Pooling layers contribute neither weights
  nor biases.
\end{solution}

% ============================================================== Question 12 ===
\question
Which of the following statements are correct (multiple answers are possible)?

\begin{choices}
  \choice In FCN-style segmentation, outputs taken at several depths of the network are upsampled
  by different factors and combined, so that fine detail from shallow layers complements the
  context of deep layers.
  \choice A label map may be transformed with bilinear interpolation provided the interpolated
  values are rounded to the nearest integer label afterwards.
  \choice Matching the mean and standard deviation of two images,
  $f(j) = (j - \mu_j)\,\sigma_i/\sigma_j + \mu_i$, matches second-order statistics and is less
  sensitive to outliers than min--max normalization.
  \choice A similarity transformation has 4 degrees of freedom in 2D and 7 in 3D.
  \choice The curl of a transformation field, $\nabla \times T$, quantifies the local volumetric
  change induced by the transformation.
\end{choices}

\begin{solution}
  Answer: (A), (C) and (D) are correct.

  (B) is wrong, and the proposed repair makes it worse rather than better. Bilinear interpolation
  between labels produces values that are not labels, and rounding them does not restore meaning:
  between a voxel of label 1 and a voxel of label 3, the interpolated value $2$ rounds to label 2
  --- possibly a completely different structure that is not even adjacent to either. Label maps are
  transformed with nearest-neighbour interpolation, full stop.

  (E) attributes the wrong quantity: the \textbf{curl} carries information about the amount of
  \emph{infinitesimal rotation}. Local volumetric change is quantified by the \textbf{Jacobian
  determinant} ($\det J > 1$ expansion, $< 1$ compression, $= 1$ none), with the divergence ---
  the trace of the Jacobian --- carrying related compression/expansion information.
\end{solution}

\end{questions}

\end{document}
